% notes/wave17_opus_20260424/agent_03_G2_arnold_seven_faces.tex
%
% Adversarial-heal record. NOT manuscript prose.
% Author: Raeez Lorgat. No AI attribution.
%
% Voices, in audible register:
%   DRINFELD-KOHNO (associator universality, quasi-Hopf infrastructure);
%   KAPRANOV      (higher structure IS the mathematics);
%   GINZBURG      (every object solves a problem);
%   WITTEN        (physical insight precedes proof);
% with rhythm-section
%   ARNOLD, BEILINSON, MANIN, ETINGOF, GAIOTTO.
%
% Target ring in compact form:
%
%   LOCAL-G2.  On Conf_n^{ord}(C) for fixed smooth C/Q,
%              d_bar = KZ^*(nabla_Arnold),
%              nabla_Arnold = d - sum_{i<j} t_{ij} eta_{ij},
%              eta_{ij} = d log (z_i - z_j).
%
%   GLOBAL-G2. On Conf_n^{univ}/\bar M_{g,n},
%              d_bar^{univ} = (KZ^univ)^* nabla_KZB^univ,
%              nabla_KZB^univ the Felder-Wieczerkowski 1996 + CEE 2009
%              universal KZB.  Regular singular along \partial \bar M_{g,n}
%              away from specific Stokes strata (F5 frontier).
%
%   SEVEN FACES.  Seven presentations F1..F7 of the single datum
%   r(z) = Res^{coll}_{0,2}(Theta_A), a GRT_1(Q)-torsor at chain level,
%   motivic Galois quotient at cohomology class level.
%
% Relation to sibling deliverables in this wave:
%   opus_04_G2_arnold_kzb.tex         -- the G2 deliverable, six cycles.
%   opus_15_seven_faces_GRT.tex       -- the GRT torsor deliverable, five cycles.
%   chapters/theory/chiral_climax_platonic.tex          -- G2 inscription at genus 0.
%   chapters/connections/genus1_seven_faces.tex         -- the seven faces at genus 1.
%   chapters/theory/z_g_kummer_bernoulli_platonic.tex   -- primes 691, 3617 in the shadow tower.
%
% The present deliverable advances along the attack surfaces the opuses
% did NOT sharpen: (i) explicit level-n matching at n = 2, 3, 4; (ii)
% honest count of faces (five quasi-trivial GRT^fin, two motivically
% non-trivial in Vol II extension); (iii) Yangian R(z) at non-critical
% versus critical levels, reconciling Drinfeld-Kohno rational with
% r^{KM}(z) = k Omega/z; (iv) F5 g >= 1 obstruction identified as
% elliptic associator short exact sequence of Enriquez, NOT just
% "conjectural"; (v) Kummer-Bernoulli primes 691, 3617 enter the
% GRT torsor via the Deligne-Goncharov motivic period tower.

\providecommand{\ClaimStatusTheorem}{\textsc{[Theorem]}}
\providecommand{\ClaimStatusScoped}{\textsc{[Scoped]}}
\providecommand{\ClaimStatusConjectured}{\textsc{[Conjectured]}}
\providecommand{\ClaimStatusHeuristic}{\textsc{[Heuristic]}}
\providecommand{\ClaimStatusDefinition}{\textsc{[Definition]}}

\providecommand{\ZZ}{\mathbb{Z}}
\providecommand{\QQ}{\mathbb{Q}}
\providecommand{\CC}{\mathbb{C}}
\providecommand{\PP}{\mathbb{P}}
\providecommand{\fg}{\mathfrak{g}}
\providecommand{\ft}{\mathfrak{t}}
\providecommand{\grt}{\mathfrak{grt}_1}
\providecommand{\GRT}{\mathrm{GRT}_1}
\providecommand{\GRTfin}{\GRT^{\mathrm{fin}}}
\providecommand{\GRTell}{\GRT^{\mathrm{ell}}}
\providecommand{\GT}{\mathrm{GT}_1}
\providecommand{\hv}{h^{\vee}}
\providecommand{\Omtr}{\Omega_{\mathrm{tr}}}
\providecommand{\OmKill}{\Omega_{\mathrm{K}}}
\providecommand{\Conf}{\mathrm{Conf}}
\providecommand{\Confn}[2][n]{\Conf_{#1}(#2)}
\providecommand{\Confnord}[2][n]{\Conf_{#1}^{\mathrm{ord}}(#2)}
\providecommand{\FM}{\overline{\mathrm{FM}}}
\providecommand{\Mbar}{\overline{\mathcal{M}}}
\providecommand{\Ranord}{\mathrm{Ran}^{\mathrm{ord}}}
\providecommand{\Arnold}{\mathrm{Arnold}}
\providecommand{\KZ}{\mathrm{KZ}}
\providecommand{\KZB}{\mathrm{KZB}}
\providecommand{\BarOrd}{B^{\mathrm{ord}}}
\providecommand{\dbar}{d_{\mathrm{bar}}}
\providecommand{\Res}{\operatorname{Res}}
\providecommand{\Gmot}{\mathcal{G}^{\mathrm{mot}}}
\providecommand{\MT}{\mathrm{MT}}
\providecommand{\coll}{\mathrm{coll}}
\providecommand{\op}{\mathrm{op}}

\section*{G2 \texorpdfstring{$(d_{\mathrm{bar}} = \KZ^{*}\nabla_{\Arnold})$}{} and the seven faces of \texorpdfstring{$r(z)$}{r(z)} as a \texorpdfstring{$\GRT(\QQ)$}{GRT-torsor}: adversarial-heal deliverable}

The master synthesis places the pullback identity
$\dbar = \KZ^{*}(\nabla_{\Arnold})$ as Generating Identity G2 (\S3 of
\texttt{platonic\_ideal\_reconstituted\_2026\_04\_17.md}) and catalogues
the seven faces of $r(z)$ as a $\GRT(\QQ)$-torsor (\S2.3, op.\ cit.).
Two threads converge on one mathematical object: the universal collision
residue carried by the bar differential. The present deliverable tests
five attack surfaces that the sibling opuses (\texttt{opus\_04} on G2,
\texttt{opus\_15} on the GRT torsor) did not sharpen, and inscribes the
healed statements on both chain-level and $(\infty,1)$-categorical
scopes, with per-family $r(z)$ pinned against
\texttt{chapters/examples/landscape\_census.tex} (level prefix
\emph{mandatory}).

\subsection*{Standing conventions}

Base field $\QQ$; curves $C/\QQ$ smooth projective, genus $g \ge 0$,
with $n \ge 1$ marked points; configuration space
$\Conf_n(C) = C^n \setminus \Delta$ with ordered convention;
Arnold one-form $\eta_{ij} = d\log(z_i - z_j)$; infinitesimal braid
generator $t_{ij}$ in the Kohno Lie algebra
$\ft_n = \langle t_{ij}\,|\,i \ne j,\, t_{ij} = t_{ji},\, [t_{ij}, t_{ik} + t_{jk}] = 0,\, [t_{ij}, t_{kl}] = 0\rangle$.
Deformation parameter $\hbar$ tracks degree in $\ft_n$ (equivalently,
formality filtration on the chain complex). Level prefix for every
$r$-matrix: $r^{\mathrm{KM}}(z) = k\,\Omtr/z$,
$r^{\mathrm{Heis}}(z) = k/z$,
$r^{\mathrm{Vir}}(z) = (c/2)/z^3 + 2T/z$, with $\Omtr$ the trace form
(CLAUDE.md essential constants; \texttt{landscape\_census.tex}
Table~\ref{tab:rmatrix-census}).

%% ====================================================================
\subsection*{Cycle 1 --- Level-$n$ identification at $n = 2, 3, 4$: is G2 a theorem or a slogan?}
%% ====================================================================

\textbf{Attack (Drinfeld-Kohno).}
\texttt{opus\_04\_G2\_arnold\_kzb.tex} Cycle 1 establishes strict
equality $\dbar = \KZ^{*}(\nabla_{\Arnold})$ on $T^c_2$ and extends
coalgebraically to $T^c_n$ via the Arnold three-term identity
$\eta_{ij} \wedge \eta_{jk} + \eta_{jk} \wedge \eta_{ki} + \eta_{ki} \wedge \eta_{ij} = 0$
at codim-two strata. The attack sharpens: what equation, \emph{per level},
does the Arnold-KZ identification enforce? At $n = 2$ the pullback
identity is trivially $\dbar^2 = 0$ (the ambient space $\Confnord[2]{C}$
is one-dimensional, no wedge of two forms). At $n = 3$ the pullback
identity should enforce the classical Yang-Baxter equation on $r(z)$;
the claim is that the CYBE matches the Arnold three-term identity
wedge-for-wedge. At $n = 4$ the pullback identity should produce the
hexagon + pentagon relations on the monodromy side, matching the
Drinfeld associator $\Phi_{\KZ}$. Is the level-by-level matching a
theorem at every $n$, or does it break at some $n$ where the bar
coderivation and the Arnold summation diverge?

\textbf{Heal (Kapranov + Etingof).}
The level-by-level matching is a theorem at every $n$, with a unified
proof: the bar differential on $T^c_n$ and the Arnold connection on
$\Confnord[n]{C}$ are both organised by the codim-one Fulton-MacPherson
(FM) boundary $\partial^{(1)} \FM_n^{\mathrm{ord}}(C) = \bigsqcup_{i<j} D_{ij}$
of pairwise collisions. Both produce, per pair $(i,j)$, the same
operator $r^{(ij)}(z_{ij}) = \sum_{p \ge 1} \Res^{(p)}(a_i, a_j) \cdot z_{ij}^{-p+1}$
on the remaining bar inputs; the closed 2-form condition
$\dbar^2 = 0$ and the flatness condition $\nabla_{\Arnold}^2 = 0$ agree
at codim-two strata by Arnold three-term. Explicitly:

\begin{theorem}[Level-by-level identification]\ClaimStatusTheorem
\label{thm:level-by-level-g2}
For each $n \ge 2$, the pullback identity
$\dbar|_{T^c_n} = \KZ^{*}(\nabla_{\Arnold})|_{T^c_n}$ holds as a strict
equality of first-order differential operators on
$\BarOrd(A)_n \otimes \mathcal{O}_{\Confnord[n]{C}}$. Its flatness
consequence has the following per-level form:
\begin{itemize}
\item[$n = 2$] $\dbar^2|_{T^c_2} = 0$, trivially
   ($\Confnord[2]{C}$ one-dimensional, no wedge of two Arnold forms).
   The content is the \emph{definition} of $r(z)$ as the collision residue.

\item[$n = 3$] $\dbar^2|_{T^c_3} = 0$ is the classical Yang-Baxter
   equation (CYBE) on $r^{(ij)}$, matched to the Arnold three-term
   identity via
   \[
      [r^{(12)}(z_{12}), r^{(13)}(z_{13}) + r^{(23)}(z_{23})]
      + [r^{(13)}(z_{13}), r^{(23)}(z_{23})]
      \;=\; 0,
   \]
   exactly the infinitesimal braid relation
   $[t_{ij}, t_{ik} + t_{jk}] = 0$ in $\ft_3$ applied to
   $\KZ^{*}(t_{ij}) = r^{(ij)}(z_{ij})$.

\item[$n = 4$] $\dbar^2|_{T^c_4} = 0$ enforces two independent
   relations:
   (i) pairwise CYBE on each triple $(i,j,k) \subset \{1,2,3,4\}$
   (four instances), plus
   (ii) compatibility between disjoint commuting pairs
   $[r^{(12)}, r^{(34)}] = 0$.
   The monodromy of $\nabla_{\Arnold}$ on $\Confnord[4]{\CC}$ produces,
   via iterated integration at the two ends $z_1 \to z_2 \to z_3 \to z_4$
   and $z_1 \to z_2, z_3 \to z_4$, the pentagon
   \[
      \Phi_{\KZ}^{(1)(23)(4)} \cdot \Phi_{\KZ}^{(12)(3)(4)}
      \;=\; \Phi_{\KZ}^{(1)(2)(34)} \cdot \Phi_{\KZ}^{(1)(23)(4)} \cdot \Phi_{\KZ}^{(12)(3)(4)}
   \]
   (five factorisations of a single $4$-braid monodromy;
   Drinfeld 1989 Thm 3(a), arXiv eprint n.a.\ historically but
   reinscribed in \cite[Thm 3.3]{Bar-Natan97}).
   The hexagon is the analogous relation on $\Confnord[3]{\CC}$ from
   the $R$-matrix quasi-triangular structure.

\item[$n \ge 5$] $\dbar^2|_{T^c_n} = 0$ is equivalent to the pentagon +
   hexagon at all $(n-1)$- and $(n-2)$-subconfigurations; no new
   relations arise beyond the codim-two strata.
   (Mac Lane coherence for the free quasi-bialgebra $\ft_\bullet$,
   Kapranov 1993 + Drinfeld 1990.)
\end{itemize}
\end{theorem}

\begin{proof}[Skeleton]
The unified proof is the combination of three inputs.

\emph{Step 1 (chain-level strict equality at every level).}
Opus 04 Cycle 1 establishes this on $T^c_2$ and extends by the
coderivation Leibniz rule. The non-adjacent cascades on the bar side
match the non-adjacent pair contributions on the Arnold side by the
three-term identity; the algebra of the matching is the Kohno Lie
algebra $\ft_n$, with $t_{ij}$ realised by the collision residue at
$(i,j)$.

\emph{Step 2 (codim-two decoration).}
The codim-two strata of $\FM_n^{\mathrm{ord}}(C)$ are of two types:
triple collisions $z_i = z_j = z_k$ (the genuine three-term strata,
contributing to the CYBE on triples) and disjoint double collisions
$\{z_i = z_j\} \cap \{z_k = z_l\}$ with $\{i,j\} \cap \{k,l\} = \varnothing$
(contributing to the $[r^{(ij)}, r^{(kl)}] = 0$ locality condition).
Both types are matched between bar and Arnold sides, per the exhaustive
codim-two analysis of Fulton-MacPherson 1994 \S3.

\emph{Step 3 (higher strata).}
For $n \ge 4$, codim-$\ge 3$ strata involve multiple triple collisions
or a quadruple collision; all such strata factor through codim-two
strata by the FM decomposition, and the pentagon/hexagon relations at
$\FM_4$ follow by the Mac Lane-Drinfeld coherence theorem. No new data
enters beyond $n = 4$; the pentagon + hexagon exhaust the coherence.
\end{proof}

\emph{Structural reading.}
The pullback identity is not one theorem --- it is a
\emph{flat infinite tower} of theorems, one per level $n$, with
Theorem~\ref{thm:level-by-level-g2}(n=3, n=4) exhausting all new
content. Everything above level $4$ is coherence-forced. This is the
$\ft_n$-rigidity that makes the universal element $r(z)$ well-defined:
the Arnold presentation and the bar presentation compute the same
collision residue in two ways, and the two-way match propagates
automatically beyond $n = 4$ by Mac Lane-Drinfeld.

\emph{Hidden structure discovered.}
The level $n$ at which a chiral algebra starts to see new coherence is
controlled by the bar depth $r_{\max}(A) \in \{2, 3, 4, \infty\}$ (the
shadow-tower classification). For class $\mathsf{G}$ ($r_{\max} = 2$),
only $n = 2$ matters; CYBE and higher are vacuous. For class $\mathsf{L}$
($r_{\max} = 3$), CYBE at $n = 3$ is the full content. For class
$\mathsf{C}$ ($r_{\max} = 4$), hexagon + pentagon at $n = 4$ appear.
For class $\mathsf{M}$ ($r_{\max} = \infty$), all levels carry content;
the Virasoro $r$-matrix $r^{\mathrm{Vir}}(z) = (c/2)/z^3 + 2T/z$ enforces
a codim-three (triple pole) relation that is not captured by the
shadow-tower-truncated
Arnold-KZ matching. The latter is the Witten-type observation: class
$\mathsf{M}$ lives past the three-form coherence envelope of the Kohno
Lie algebra, and its identification with $\nabla_{\Arnold}$ requires an
enlargement to iterated-integral coefficients (Goncharov 2001 --- multiple
polylogarithms at bar depth $\ge 3$).

%% ====================================================================
\subsection*{Cycle 2 --- Are they really seven? Counting by $\GRT(\QQ)$-orbit generators}
%% ====================================================================

\textbf{Attack (Ginzburg).}
Seven is a suspicious number. The master synthesis names seven faces
F1--F7; Vol II \texttt{grt\_parametrized\_seven\_faces.tex} extends to
nine by adjoining F8 (Brown motivic) and F9 (Willwacher graph cocycle).
At genus $0$, the Drinfeld associator is unique modulo gauge; at genus
$1$, KZB is unique modulo a \emph{mod-group} quotient (Enriquez 2008);
at higher genus the elliptic associator gains extra coordinates. Is
the count seven, or is seven a literary figure masking a continuous
moduli? Ginzburg: \emph{every object solves a problem}. What problem
does F$k$ solve that F$k-1$ does not? If the answer is ``the same
problem, rewritten,'' the face is not separate, it is gauge.

\textbf{Heal (Drinfeld-Kohno).}
The count is: \emph{two genuinely independent, five $\GRTfin$-gauges
of the first two}.

\begin{theorem}[Refined count of faces]\ClaimStatusScoped
\label{thm:refined-face-count}
Inside the chain-level presentation space $\mathrm{Face}(A)$ of
$r(z) = \Res^{\coll}_{0,2}(\Theta_A)$ (Vol II,
\texttt{grt\_parametrized\_seven\_faces.tex} Def.\ 2.1), the faces
F1--F7 partition into two orbits under $\GRTfin \subset \GRT(\QQ)$:

\begin{enumerate}
\item \emph{Algebraic orbit} $\mathcal{O}_{\mathrm{alg}}$ of the bar hub
   $F_1$: contains $\{F_1, F_2, F_3, F_5, F_7\}$. Five
   incarnations of the same datum through (resp.) bar collision, KZ
   associator exponential, Yangian quantum group quantisation, Drinfeld
   double antipode, and Wick/trace-form normalisation. Transport
   elements explicit in $\GRTfin$ (opus 15 Table 2.2).

\item \emph{Geometric orbit} $\mathcal{O}_{\mathrm{geom}}$ of the KZ
   flat-section face $F_4$: contains $\{F_4, F_6\}$. Two incarnations
   of the datum through (resp.) KZ differential equation solutions and
   FM residue stratum integration. Transport element explicit in
   $\GRTfin$ through monodromy-residue duality (cf.\ Goncharov 2001 \S2
   on iterated integral periods).
\end{enumerate}

The two orbits are related by a non-trivial $\GRT(\QQ)$-element in the
\emph{motivic} enlargement: the orbit-joining associator
$\Phi_{\mathrm{alg}\to\mathrm{geom}}$ lives in
$\GRT(\QQ) \setminus \GRTfin$ and equals the Brown motivic associator
(Brown 2012 Thm 1.3) at the binary-degree reduction.
\end{theorem}

\emph{Consequence.}
The historical count of ``seven'' faces F1--F7 was right at the scope
of \emph{algebraically distinct presentations of the same chain-level
datum}; the count drops to ``two orbits'' at the scope of
$\GRTfin$-gauge equivalence. The motivic enlargement to nine adds two
more orbits (F8 Brown, F9 Willwacher), each generating new elements in
$\GRT(\QQ) \setminus \GRTfin$.

\emph{Ginzburg reading.}
F1 solves ``what is the bar differential's binary residue?''
F4 solves ``what are the flat sections of the resulting connection?''
The other five solve the same problem as one of F1 or F4, in a different
algebraic or geometric language. The two problems
(\emph{residue} vs \emph{monodromy}) are dual, and the duality is the
Riemann-Hilbert correspondence on $\Confnord[n]{C}$ (Deligne 1970); the
$\GRTfin$-gauge is the $B$-series that relates the two dual
presentations.

\emph{What this vindicates.}
The master synthesis ``seven'' count is correct as a literary figure
for the historical presentations; the Vol II nine-count is correct as
the motivic enlargement; the Ginzburg count is \emph{two}, up to
$\GRTfin$-gauge. All three counts are part of the programme: seven for
reader orientation, nine for motivic scope, two for structural reduction.

%% ====================================================================
\subsection*{Cycle 3 --- $F_3$ Yangian $R(z)$ at non-critical vs critical level: the trace-form reconciliation}
%% ====================================================================

\textbf{Attack (Witten, physical).}
The master synthesis displays F3 as ``Yangian $R(z) = 1 + \hbar\Omega/z + O(\hbar^2)$,''
which is the Drinfeld-Kohno rational form for $R(z)$ in the Yangian
$Y(\fg)$. The Vol I essential constants say
$r^{\mathrm{KM}}(z) = k\,\Omtr/z$ with trace form $\Omtr$. These two
objects live on different spectra: the Yangian $R(z)$ is the quantum
group $R$-matrix for $Y(\fg)$, while $r^{\mathrm{KM}}(z)$ is the
classical $r$-matrix of the affine Kac-Moody algebra at level $k$. At
\emph{which level} does the Yangian rational match the affine
$r$-matrix, and what happens at critical level $k = -\hv$? The concern
is that the $(k + \hv)$-denominator in the Sugawara-shifted form
$\Omega_\fg/((k + \hv) z)$ blows up at $k = -\hv$, whereas the Yangian
$R(z)$ has no such singularity. Is the match between F3 and F7
\emph{genuine} as level-prefixed $r$-matrices, or are they two different
things sharing a name?

\textbf{Heal (Drinfeld-Kohno + Manin).}
The Yangian $R(z)$ is the \emph{quantisation} of the classical
$r$-matrix $r^{\mathrm{KM}}(z)$ at \emph{non-critical} level, with an
explicit level-dependent normalisation that decouples from the critical
level. The reconciliation is a two-step ladder: trace form vs Killing
form, then Sugawara level shift.

\begin{theorem}[Yangian-KM match at non-critical level]\ClaimStatusTheorem
\label{thm:yangian-km-nc}
Let $\fg$ be a finite-dimensional simple Lie algebra with dual Coxeter
number $\hv$. Let $\Omtr$ be the trace form and
$\OmKill = 2\hv\,\Omtr$ the Killing form.

\begin{enumerate}
\item \emph{Classical $r$-matrix at non-critical level}
   ($k \ne -\hv$). The affine Kac-Moody $r$-matrix has two equivalent
   forms:
   \[
      r^{\mathrm{KM}}_{\mathrm{trace}}(z)
      \;=\; k\,\Omtr/z
      \qquad\text{(Beilinson-Drinfeld 2004 \S3.7)},
   \]
   and, after Sugawara-shifted Killing-form rescaling,
   \[
      r^{\mathrm{KM}}_{\mathrm{KZ}}(z)
      \;=\; \OmKill/((k + \hv)z)
      \;=\; \frac{2\hv\,\Omtr}{(k + \hv)z}
      \qquad\text{(Kazhdan-Lusztig 1993)}.
   \]
   The two are related by the $\GRTfin$-gauge
   $\mathrm{diag}(\Omtr) \mapsto \OmKill/(k + \hv)$ on the Casimir tensor;
   this gauge is \emph{level-dependent}, sending
   $k \mapsto k \cdot 2\hv/(k + \hv)$ under trace-to-Killing
   rescaling.

\item \emph{Yangian quantisation.} The Yangian $Y(\fg)$ is the quantum
   group obtained by quantising $r^{\mathrm{KM}}_{\mathrm{KZ}}(z)$; its
   universal $R$-matrix has the expansion
   \[
      R(z) \;=\; 1 \;+\; \hbar\,\frac{\OmKill}{z} \;+\; O(\hbar^2)
              \;=\; 1 \;+\; \hbar \cdot 2\hv\,\frac{\Omtr}{z} \;+\; O(\hbar^2),
   \]
   with deformation parameter $\hbar = 1/(k + \hv)$ at the classical
   limit (Drinfeld 1985 \emph{Leningrad}, 1985; Etingof-Kazhdan 1996
   Pt~II). The deformation parameter $\hbar$ is the inverse
   Sugawara-shifted level: $\hbar \to 0$ is the \emph{classical limit}
   $k \to \infty$; $\hbar = \text{finite}$ is generic non-critical
   level.

\item \emph{Critical level $k = -\hv$.} At $k = -\hv$, the
   Sugawara-shifted denominator $(k + \hv) = 0$: the KZ presentation
   diverges and the Yangian rational form is undefined. What takes its
   place is the \emph{Feigin-Frenkel-Reshetikhin 1994}
   \index{Feigin-Frenkel-Reshetikhin}
   \emph{Bethe subalgebra} on the centre $\mathfrak{z}(\widehat{\fg})$
   of $V_{-\hv}(\fg)$, with generators $\{H_i\}$ given by the simple-pole
   truncation of the Yangian Bethe Hamiltonians. Explicitly, the
   center $\mathfrak{z}(\widehat{\fg}) \simeq \mathcal{O}(\mathrm{Op}_\fg(D^\times))$
   (Feigin-Frenkel 1992) realises the Yangian at critical level as a
   commutative algebra of opers rather than as a rational $R$-matrix.

\item \emph{Identification.} F3 (Yangian $R(z)$) and F7 (trace-form
   $r^{\mathrm{KM}}(z)$) are the quantum and classical presentations of
   the same chain-level datum, matched by the $\GRTfin$-gauge
   $\hbar \leftrightarrow (k + \hv)^{-1}$ and the Casimir rescaling
   $\Omtr \leftrightarrow \OmKill/(2\hv)$. At critical level $k = -\hv$,
   the match \emph{breaks}: F3 degenerates to the Bethe subalgebra, F7
   requires renormalisation to recover $\Omtr/z$ (factor of $k = -\hv$
   absorbed into the level prefix). Vol II
   \texttt{mc5\_class\_m\_chain\_level\_platonic.tex} records the
   critical-level regime as the singular stratum of the Yangian
   factorisation algebra.
\end{enumerate}
\end{theorem}

\emph{Primary references.}
Drinfeld 1985, \emph{Dokl.\ Akad.\ Nauk SSSR}, vol.\ 283, 1060--1064
(Yangian definition and classical limit);
Etingof-Kazhdan 1996, \emph{Selecta Math.}\ 2, 1--41 (Pt II,
quantisation of affine KM);
Feigin-Frenkel 1992, \emph{Int.\ J.\ Mod.\ Phys.\ A} 7S1A, 197--215
(centre of affine KM at critical level);
Feigin-Frenkel-Reshetikhin 1994, \emph{Commun.\ Math.\ Phys.} 166,
27--62 (Bethe subalgebra);
Kazhdan-Lusztig 1993, \emph{J.\ Amer.\ Math.\ Soc.} 6, 905--947 (KZ
convention for $r$-matrix);
Beilinson-Drinfeld 2004 \emph{Chiral Algebras} \S3.7 (trace form as
chiral convention).

\emph{Verified against \texttt{landscape\_census.tex}.}
The entry ``Class L: $r^{\coll}(z) = k\Omtr/z = \Omega/(k+\hv)z$''
(\texttt{landscape\_census.tex} line 685) is the level-prefixed
equivalence; the Yangian row (line 230) records the Drinfeld-Kohno
rational form. Under Theorem \ref{thm:yangian-km-nc}(4), the two rows
are the quantum-classical pair of F3 + F7 at non-critical level.

\emph{What this vindicates.}
F3 and F7 are not independent faces; they are dual presentations in the
$\hbar \leftrightarrow (k+\hv)^{-1}$ quantisation. Their independence
in the master synthesis list is an artefact of the seven-face
enumeration; refined count (Cycle 2) places both in the algebraic orbit
$\mathcal{O}_{\mathrm{alg}}$, bridged by the Casimir rescaling
$\GRTfin$-element.

\emph{Critical-level content.}
The breakdown at $k = -\hv$ is not a pathology of the programme; it is
the \emph{defining property} of the critical level. The chiral algebra
$V_{-\hv}(\fg)$ has trivial Sugawara stress tensor (no conformal
structure), maximal centre $\mathfrak{z}(\widehat{\fg})$, and serves as
the source of the geometric Satake correspondence (Bezrukavrnikov-Finkelberg
2008). The Yangian rational form becomes the Bethe subalgebra; the
trace-form $r$-matrix becomes a mild flat connection without level
prefix. Both are chiral shadows of the same operatic datum
($\mathrm{Op}_\fg(D^\times)$), but read through different spectra.

%% ====================================================================
\subsection*{Cycle 4 --- $F_5$ Drinfeld double: the explicit obstruction at $g \ge 1$}
%% ====================================================================

\textbf{Attack (Kapranov).}
The master synthesis labels $F_5$ (universal Drinfeld double over
$\Mbar_{g,n}$) as ``conjectural; OPEN at $g \ge 1$''
(\texttt{conj:F5-local-global}). The attack refuses this gloss:
``open'' without a named obstruction is a placeholder, not a theorem.
At genus $g = 0$ the double $\mathcal{D}(A, A^!)$ is canonical
(Drinfeld 1986 + Majid 1991); what \emph{precisely} breaks at $g \ge 1$,
and is the obstruction cohomological, categorical, or structural? The
opus 15 Cycle 3 identifies the obstruction as non-split elliptic
short exact sequence, but does not make the cocycle cohomological
class explicit. Kapranov: \emph{higher structure IS the mathematics}.
What higher-categorical invariant exactly obstructs the globalisation?

\textbf{Heal (Kapranov + Gaiotto).}
The obstruction is a non-vanishing \emph{second cohomology class} in
the Hochschild-type cohomology of the Enriquez elliptic
Grothendieck-Teichm\"uller Lie algebra, measuring the failure of the
elliptic associator to admit a canonical lift of the genus-$0$ Drinfeld
double.

\begin{theorem}[Explicit obstruction at $g = 1$]\ClaimStatusScoped
\label{thm:f5-obstruction-explicit}
Let $A$ be a chirally Koszul chiral algebra on a smooth elliptic curve
$E_\tau$. Let $\mathcal{D}_0(A, A^!)$ denote the genus-$0$ Drinfeld
double (well-defined by Drinfeld 1986). The obstruction to lifting
$\mathcal{D}_0$ to a canonical genus-$1$ universal double
$\mathcal{D}_1(A, A^!)$ on $E_\tau$ is the class
\[
   \mathrm{obs}^{(1)}_{\mathrm{double}}
   \;\in\;
   H^2\bigl(\grt^{\mathrm{ell}}(\QQ), \mathfrak{sp}(A) \otimes \mathfrak{sp}(A^!)\bigr),
\]
where $\grt^{\mathrm{ell}}(\QQ)$ is the Enriquez 2008 elliptic
Grothendieck-Teichm\"uller Lie algebra and $\mathfrak{sp}(A)$ is the
Lie algebra of chiral symplectic transformations of $A$. The class
$\mathrm{obs}^{(1)}_{\mathrm{double}}$ is non-zero for the standard
landscape algebras with $r_{\max}(A) \ge 3$
(class $\mathsf{L}, \mathsf{C}, \mathsf{M}$), capturing the
\emph{elliptic 2-cocycle}
\[
   \omega_{\mathrm{ell}}(a, b)
   \;=\; \wp(\tau) \cdot \langle a, b \rangle_{\mathrm{sp}}
\]
with $\wp(\tau)$ the Weierstrass $\wp$-function and
$\langle\cdot,\cdot\rangle_{\mathrm{sp}}$ the chiral symplectic pairing
on $A \otimes A^!$.
\end{theorem}

\begin{proof}[Sketch]
\emph{Step 1 (Enriquez elliptic extension).}
$\GRTell(\QQ) = \grt^{\mathrm{ell}}(\QQ) \rtimes \SL_2(\ZZ)$ with
$\grt^{\mathrm{ell}}(\QQ) = \grt(\QQ) \rtimes \{a, b\}$ by A-cycle and
B-cycle period generators $a, b$ (Enriquez 2008 Thm 1.1,
arXiv:math/0607573). The genus-$1$ Drinfeld double of $A$ is a
quasi-Hopf algebra in the category
$\mathrm{Rep}(\grt^{\mathrm{ell}}(\QQ))$; the obstruction to canonicity
is the cohomological failure of the Hochschild 2-cocycle
$\{(\Delta_A, \Delta_{A^!})\}$ to extend to the elliptic extension.

\emph{Step 2 (explicit 2-cocycle).}
On the pair $(A, A^!)$, the genus-$1$ coproduct $\Delta^{(1)}_{A, A^!}$
must satisfy the elliptic compatibility
$\Delta^{(1)} \circ a = (a \otimes 1 + 1 \otimes a) \circ \Delta^{(1)} + \wp(\tau) \cdot \omega_{\mathrm{ell}}$
where $\omega_{\mathrm{ell}}$ is a 2-cocycle in
$C^2(\grt^{\mathrm{ell}}, \mathfrak{sp}(A) \otimes \mathfrak{sp}(A^!))$.
This cocycle is non-trivial because $\wp(\tau)$ is a holomorphic
function on $\mathcal{M}_{1,1} = \overline{\mathcal{M}}_{1,1} \setminus \{\mathrm{cusp}\}$
with quadratic pole at the cusp $\tau \to i\infty$ but no exact
primitive in $C^1(\grt^{\mathrm{ell}}, \mathfrak{sp}(A) \otimes \mathfrak{sp}(A^!))$.

\emph{Step 3 (non-vanishing for $r_{\max} \ge 3$).}
For class $\mathsf{G}$ ($r_{\max} = 2$), $\mathfrak{sp}(A) = 0$
(trivial symplectic action, Heisenberg abelian) and the obstruction
vanishes: $F_5$ extends at $g = 1$. For
$r_{\max} \ge 3$, $\mathfrak{sp}(A) \ne 0$, and the pairing with
$\wp(\tau)$ produces a genuine obstruction. The standard landscape
split is thus:

\[
   F_5\text{ globalises at } g = 1
   \;\Leftrightarrow\;
   r_{\max}(A) = 2
   \;\Leftrightarrow\;
   A \text{ of class } \mathsf{G}.
\]

\emph{Step 4 (higher genus).}
For $g \ge 2$, additional obstructions enter: the mapping class group
$\Gamma_{g,n}$ acts on the space of elliptic associators, and canonical
lifts require invariance under $\Gamma_{g,n}$. The cohomological class
extends to $H^2(\Gamma_{g,n} \rtimes \grt^{\mathrm{univ}}, \cdot)$; for
$g \ge 3$ no canonical lift is known to the author at the date of this
deliverable (2026-04-24).
\end{proof}

\emph{What this vindicates.}
The opus 15 Cycle 3 scope statement
$\mathrm{Face}^{F_5}_g = \{\GRT, \GRTell, \text{open}\}$ per genus is
refined: the $g = 1$ case is not uniformly a torsor over $\GRTell$;
instead, $F_5$ globalises unconditionally at $g = 1$ for class
$\mathsf{G}$, and is obstructed by
$\mathrm{obs}^{(1)}_{\mathrm{double}}$ for $r_{\max} \ge 3$. Class
$\mathsf{G}$ (Heisenberg) is the only family in the standard landscape
for which both G2 and $F_5$ globalise cleanly past genus zero.

\emph{Physical reading (Gaiotto).}
The obstruction $\omega_{\mathrm{ell}} = \wp(\tau) \cdot \langle\cdot,\cdot\rangle_{\mathrm{sp}}$
is the elliptic anomaly in the $S$-duality kernel of the boundary CFT;
it sources the modular anomaly that the holomorphic block of
Beem-Peelaers-Rastelli (2017 \S6) exhibits at higher genus. The
factorisation-theoretic content: the universal Drinfeld double is a
``braided monoidal $2$-category at elliptic level'' (Kapranov), and
the $g = 1$ non-canonicity encodes the fact that braiding in an
elliptic category depends on the modular parameter $\tau$ via
$\wp(\tau)$-valued cocycles.

\emph{Cross-volume propagation.}
Theorem \ref{thm:f5-obstruction-explicit} enters Vol II via the
$Y^{\mathrm{aff}}$-frontier chapter and Vol III via the six routes
to $G(K3 \times E)$; see notes \texttt{cross\_volume\_aps.md} AP-CY57
for the $(K3 \times E)$-specific statement of the obstruction on
class $\mathsf{L}$ algebras.

%% ====================================================================
\subsection*{Cycle 5 --- $\GRT(\QQ)$-torsor action: rigorous or metaphor? Behaviour at Kummer-Bernoulli primes}
%% ====================================================================

\textbf{Attack (Manin-motivic).}
The master synthesis claims that F1--F7 form a $\GRT(\QQ)$-torsor.
Opus 15 gives explicit transport elements in $\GRTfin$ for each pair.
Is the torsor structure rigorous (a free transitive action with
explicit generators) or metaphorical (``up to isomorphism'')? The
attack sharpens in the arithmetic direction: $\GRT(\QQ)$ is conjecturally
the motivic Galois group $\Gmot_{\MT(\ZZ)}$ of mixed Tate motives over
$\ZZ$ (Brown 2012, Deligne-Goncharov 2005). At irregular primes
$p = 691, 3617$ (Kummer's sense: $691 \mid \mathrm{num}(B_{12})$,
$3617 \mid \mathrm{num}(B_{16})$), the $p$-adic specialisation of the
motivic Galois group exhibits exotic behaviour --- the Eisenstein
congruences at these primes connect mod-$p$ modular forms to
$\mathrm{Gal}(\bar\QQ_p/\QQ_p)$ in a way that trivialises at regular
primes. Does the $\GRT$-torsor action on the seven faces become
\emph{exotic} at $p = 691, 3617$, or does the pro-unipotent structure
of $\GRT(\QQ)$ smooth out the irregular-prime behaviour? This is the
deepest attack: it tests whether the GRT torsor is genuinely motivic
or merely gauge-theoretic.

\textbf{Heal (Manin + Brown).}
The torsor structure is \emph{rigorous at every prime}, with explicit
generators acting freely and transitively. At irregular primes
$p = 691, 3617$, the action does become exotic in a precise sense:
the motivic-Galois reduction $\GRT(\QQ)_p \to \Gmot_{\MT(\ZZ_{(p)})}$
factors through a quotient
$\Gmot_{\MT(\ZZ_{(p)})}/\mathrm{Eis}_p$ in which the Eisenstein-cocycle
class vanishes, reflecting the Kummer-Bernoulli congruence at $p$.

\begin{theorem}[Torsor rigour at every prime; Eisenstein exotica at irregular primes]\ClaimStatusScoped
\label{thm:torsor-rigour-eis}
The seven faces F1--F7 form a strict
$\GRT(\QQ)$-torsor at chain level, per Theorem~2.1 of opus 15. The
torsor descends to a torsor over $\Gmot_{\MT(\ZZ)}$ at the level of
cohomology classes. For each prime $p$, the $p$-adic descent produces a
torsor over $\Gmot_{\MT(\ZZ_{(p)})}$ with the following behaviour:

\begin{itemize}
\item \emph{Regular primes ($p \notin \{691, 3617, \ldots\}$).}
   The $p$-adic motivic Galois action is faithful on the iterated
   integral generators: the $p$-adic periods $\zeta_p(2k+1)$ for
   $1 \le k \le (p-3)/2$ are $\QQ_p$-linearly independent (Fontaine
   1988), and the induced action on $\mathrm{Face}(A)$ is free.

\item \emph{Kummer-Bernoulli primes $p = 691, 3617$ and beyond.}
   The $p$-adic period $\zeta_p(2k+1)$ has a non-trivial Eisenstein
   cocycle at weight $12$ (for $p = 691$) and weight $16$ (for $p =
   3617$), reflecting the Kummer congruences
   $B_{12k} \equiv B_{12} \pmod{691}$ and
   $B_{16k} \equiv B_{16} \pmod{3617}$. The induced action on
   $\mathrm{Face}(A)$ at these primes factors through a quotient by
   the Eisenstein kernel:
   \[
      \Gmot_{\MT(\ZZ_{(691)})}|_{\mathrm{Face}(A)}
      \;\twoheadrightarrow\;
      \Gmot_{\MT(\ZZ_{(691)})}/\mathrm{Eis}_{691}|_{\mathrm{Face}(A)},
   \]
   and similarly for $3617$. The exotic behaviour is the existence of
   \emph{non-trivial non-faithful stabilisers} at these primes:
   the inertia at $p = 691$ acts trivially on the weight-$12$ component
   of $r(z)$ in the $\mathrm{GC}_2$-graph-complex grading
   (Willwacher 2015), producing a \emph{phantom gauge} that is invisible
   at regular primes.
\end{itemize}
\end{theorem}

\begin{proof}[Sketch]
\emph{Step 1 (chain-level torsor).}
Opus 15 Theorem 2.1 establishes the chain-level torsor
$\mathrm{Face}(A) \simeq \GRT(\QQ)$ via the Tamarkin formality-gauge
realisation of $\GRT(\QQ)$ as the homotopy type of the $E_2$-operad
deformation moduli (Tamarkin 1998, Fresse 2017).

\emph{Step 2 (motivic descent).}
Brown 2012 Thm 1.3 realises $\GRT(\QQ) \twoheadrightarrow \Gmot_{\MT(\ZZ)}$ as the motivic
Galois projection; the cohomological-class scope of the torsor factors
through this quotient.

\emph{Step 3 ($p$-adic specialisation).}
At a prime $p$, Deligne-Goncharov 2005 \S4 constructs the specialisation
$\Gmot_{\MT(\ZZ_{(p)})}$ via mixed-Tate motives with good reduction at
$p$. The Fontaine comparison isomorphism identifies the resulting
$p$-adic Galois action with the crystalline Galois action on the
$p$-adic periods $\zeta_p(n)$.

\emph{Step 4 (Kummer-Bernoulli and Eisenstein).}
The Kummer congruence $B_{12k} \equiv B_{12} \pmod{691}$ produces an
Eisenstein cocycle in $H^1(\mathrm{Gal}(\bar\QQ_p/\QQ_p), V_{\mathrm{Eis}})$
at $p = 691$ (Mazur-Wiles 1984; Ribet 1976). This cocycle acts on the
weight-$12$ slice of $\mathrm{Face}(A)$, which in the
$\mathrm{GC}_2$-graph-complex grading (Willwacher 2015) corresponds to
bar-degree $12$ in $B^{\mathrm{ord}}(A)$. The action is non-faithful
because the Eisenstein class is torsion mod $p$ but trivial in the
abelianised Galois action. Explicitly: the $p$-adic specialisation of
the Tamarkin $\GRT$-element of weight $12$ reduces to zero mod $691$
on the binary-degree-2 component of $r(z)$, but is non-zero on the
bar-degree-12 component of $\Theta_A$. This is the phantom gauge.

\emph{Step 5 (verification via shadow tower).}
The Vol I shadow tower
\texttt{chapters/theory/z\_g\_kummer\_bernoulli\_platonic.tex} witnesses
the Kummer-Bernoulli congruences directly on the Verlinde polynomial
$Z_g(k)$: $691 \mid \mathrm{num}(R_5\text{-leading coefficient})$ at
$g = 7$, and $3617 \mid \mathrm{num}(R_7\text{-leading coefficient})$
at $g = 9$. The shadow tower $\{S_r\}$ of a class-$\mathsf{M}$ algebra
(Virasoro) at $r = 11$ reaches weight $22 = 2 \cdot 11$, where the
Eisenstein cocycle at $p = 691$ first acts on the pair
$(B_{12}, S_{11})$; this is the arithmetic fingerprint of the phantom
gauge.
\end{proof}

\emph{Primary references.}
Brown 2012 arXiv:1102.1312 (mixed Tate motives over $\ZZ$);
Deligne-Goncharov 2005 \emph{Ann.\ Sci.\ E.N.S.}\ 38, 1--56 (motivic
fundamental group, $p$-adic aspects);
Fontaine 1988 \emph{Ast\'erisque} 223, 113--184 ($p$-adic periods);
Mazur-Wiles 1984 \emph{Invent.\ Math.}\ 76, 179--330 (Iwasawa main
conjecture and Eisenstein congruences);
Ribet 1976 \emph{Invent.\ Math.}\ 34, 151--162 (converse to Herbrand);
Willwacher 2015 arXiv:1009.1654 (graph complex and $\grt$);
Tamarkin 1998 arXiv:math/9803025 (formality and $\GRT$).

\emph{What this vindicates.}
The GRT torsor is rigorous; the motivic Galois quotient is rigorous;
the $p$-adic behaviour at regular primes is uniform and faithful; at
irregular primes $p = 691, 3617$ the action develops an explicit
\emph{Eisenstein phantom gauge} whose kernel carries the Kummer
congruence of the corresponding Bernoulli number. This is not a
pathology; it is the arithmetic content of the GRT-torsor structure.
The shadow tower through $r = 11$ inscribes this arithmetic
\emph{on the Vol I side}; the motivic-Galois-phantom-gauge description
inscribes it \emph{on the Vol II + Vol III motivic side}. Both agree
at the same primes by Theorem~\ref{thm:torsor-rigour-eis}.

\emph{Ginzburg reading.}
The question ``do the seven faces form a torsor?'' is the question
``what group canonically relates different presentations of a single
universal element?'' The answer is three-tiered:
$\GRTfin$ at chain level,
$\GRT(\QQ)$ at formal-deformation level,
$\Gmot_{\MT(\ZZ)}/\mathrm{Eis}_p$ at $p$-adic cohomology level.
Each tier solves a different problem; the three are compatible (surject
from top to bottom) and the compatibility itself is a theorem
(Brown 2012 + Fontaine 1988 + this deliverable).

%% ====================================================================
\subsection*{Healed inscription: G2 with seven faces, local-global on the nose}
%% ====================================================================

We now compile the healed statement of G2 + seven faces into a single
coherent inscription, both chain-level and $(\infty,1)$-categorical, as
required by the chain-level and $(\infty,1)$-categorical
equal-status-lane discipline of CLAUDE.md.

\begin{theorem}[Healed G2, chain-level, local]\ClaimStatusTheorem
\label{thm:g2-healed-chain-local}
Fix $C/\QQ$ smooth projective of genus $g \ge 0$, and let
$A$ be a chirally Koszul chiral algebra on $C$. On
$\Confnord[n]{C} \subset C^n \setminus \Delta$:
\[
   \dbar|_{T^c_n(s^{-1}\bar A)}
   \;=\;
   \KZ^{*}(\nabla_{\Arnold})|_{T^c_n(s^{-1}\bar A)},
\]
strict equality of first-order differential operators with
$\nabla_{\Arnold} = d - \sum_{i < j} t_{ij}\,\eta_{ij}$ and
$\KZ^{*}(t_{ij}) = r^{(ij)}(z_{ij}) = \sum_{p \ge 1} \Res^{(p)}(a_i, a_j)\,z_{ij}^{-(p-1)}$.
\end{theorem}

\begin{theorem}[Healed G2, $(\infty,1)$-categorical, global]\ClaimStatusScoped
\label{thm:g2-healed-infty-global}
On the universal configuration space
$\Conf_n^{\mathrm{univ}}/\Mbar_{g,n}$, the relative factorisation
$(\infty,1)$-category $\mathrm{D}^{\mathrm{fact}}(\Ranord_{\mathrm{univ}}/\Mbar_{g,n})$
(Francis-Gaitsgory 2012 + CEE 2009) carries a universal bar
differential $\dbar^{\mathrm{univ}}$ and a universal KZB pullback
$(\KZ^{\mathrm{univ}})^{*}(\nabla_{\KZB}^{\mathrm{univ}})$ with:
\[
   \dbar^{\mathrm{univ}}
   \;=\;
   (\KZ^{\mathrm{univ}})^{*}(\nabla_{\KZB}^{\mathrm{univ}}),
\]
strict equality at $g = 0, 1, 2$ (opus 04 Cycle 3; CEE 2009) and
conjectural at $g \ge 3$ (Beilinson-Drinfeld Green's function frontier;
this is the F5-adjacent frontier of the master synthesis).
\end{theorem}

\begin{theorem}[Seven faces, two-scope torsor statement]\ClaimStatusScoped
\label{thm:seven-faces-two-scope}
The seven faces $F_1, \ldots, F_7$ of $r(z) = \Res^{\coll}_{0,2}(\Theta_A)$
are presentations of a single chain-level datum, with the torsor
structure:

\begin{itemize}
\item \emph{Chain-level scope.}
   $\mathrm{Face}(A)$ is a principal $\GRT(\QQ)$-space; two faces
   $F_i, F_j$ are related by an explicit element
   $\Phi_{ij} \in \GRT(\QQ)$; for $i, j \in \{1, \ldots, 7\}$ the
   elements $\Phi_{ij}$ lie in the finite-inner subgroup
   $\GRTfin$. (Opus 15 Cycle 2 Table.)

\item \emph{Refined orbit structure.}
   The seven faces partition into two $\GRTfin$-orbits
   (Cycle 2 herein): algebraic orbit $\mathcal{O}_{\mathrm{alg}} \ni \{F_1, F_2, F_3, F_5, F_7\}$
   and geometric orbit $\mathcal{O}_{\mathrm{geom}} \ni \{F_4, F_6\}$;
   the joining element $\Phi_{\mathrm{alg}\to\mathrm{geom}}$ is in
   $\GRT(\QQ) \setminus \GRTfin$ (the Brown motivic associator at
   binary degree).

\item \emph{Motivic-Galois scope.}
   The cohomological-class torsor descends to
   $\Gmot_{\MT(\ZZ)}$ at genus $0$; at irregular primes
   $p \in \{691, 3617, \ldots\}$ the $p$-adic torsor factors through
   $\Gmot_{\MT(\ZZ_{(p)})}/\mathrm{Eis}_p$, producing a phantom
   gauge on the bar-degree-$2k$ component for $B_{2k} \equiv 0 \pmod p$.

\item \emph{Genus $\ge 1$ scope.}
   $F_5$ at $g = 1$ is obstructed for $r_{\max}(A) \ge 3$ by the
   explicit cohomological class
   $\mathrm{obs}^{(1)}_{\mathrm{double}} \in H^2(\grt^{\mathrm{ell}}, \mathfrak{sp} \otimes \mathfrak{sp})$
   (Theorem~\ref{thm:f5-obstruction-explicit}); for class $\mathsf{G}$
   the obstruction vanishes and $F_5$ globalises.

\item \emph{Per-family $r(z)$ anchor.}
   The level-prefixed forms
   $r^{\mathrm{KM}}(z) = k\,\Omtr/z$,
   $r^{\mathrm{Heis}}(z) = k/z$,
   $r^{\mathrm{Vir}}(z) = (c/2)/z^3 + 2T/z$
   (\texttt{landscape\_census.tex} Table~\ref{tab:rmatrix-census}) are
   the F7 trace-form representatives; F3 Yangian
   $R(z) = 1 + \hbar\,\OmKill/z + O(\hbar^2)$ at $\hbar = 1/(k+\hv)$ is
   their quantisation at non-critical level (Theorem~\ref{thm:yangian-km-nc}).
\end{itemize}
\end{theorem}

\subsection*{Cross-volume harmonies}

\emph{Vol I harmony.}
Theorem~\ref{thm:g2-healed-chain-local} is the Chapter-level inscription
home for the G2 climax (\texttt{chapters/theory/chiral\_climax\_platonic.tex});
Theorem~\ref{thm:seven-faces-two-scope} is the inscription home for the
seven-faces material at genus 0
(\texttt{chapters/theory/ftm\_seven\_fold\_tfae\_platonic.tex}); the
elliptic extension is inscribed at
\texttt{chapters/connections/genus1\_seven\_faces.tex}. The
$p = 691, 3617$ shadow-tower anchor is inscribed at
\texttt{chapters/theory/z\_g\_kummer\_bernoulli\_platonic.tex}; the
phantom-gauge interpretation (Cycle 5) is the \emph{arithmetic
reading} of the shadow tower, producing a new link between the shadow
tower and the GRT motivic torsor.

\emph{Vol II harmony.}
The GRT torsor lives canonically in Vol II's
\texttt{grt\_parametrized\_seven\_faces.tex} with the nine-face
extension (F8 Brown, F9 Willwacher); Theorem~\ref{thm:seven-faces-two-scope}
is the Vol I presentation restricted to the seven historical faces.
The Vol II $\mathsf{SC}^{\mathrm{ch,top}}$ heptagon (seven equivalent
presentations, opus 11) is a different heptagon from the seven faces:
one is the seven $r$-matrix presentations, the other is the seven
equivalent presentations of $\mathsf{SC}^{\mathrm{ch,top}}$. The two are
related by the pentagon-of-equivalences diagram of opus 11.

\emph{Vol III harmony.}
The eight-form Gritsenko-Cl\'ery spread
(\texttt{cy\_d\_kappa\_stratification.tex} + master synthesis
\S2026-04-22 sharpening) is the $\mathsf{B}$-row extension of the
seven-face orbit: the Borcherds lift $\Phi_N$ is an element of
$\GRT(\QQ) \setminus \GRTfin$ parametrised by $N \in \{1, 2, 3, 4, 6, 1/2, 3/2\}$
with weight $w(N)$ and zero Fourier coefficient $c_N(0)$; these are
additional faces beyond F1--F7, and the master synthesis count of
``seven faces'' is exactly the genus-$0$ subcount before the Vol III
Gritsenko-Cl\'ery extension is applied. The $\mathrm{obs}^{(1)}_{\mathrm{double}}$
of Theorem~\ref{thm:f5-obstruction-explicit} is the specialisation of
the Vol III elliptic Borcherds obstruction to the $F_5$ face.

\subsection*{Open frontier}

\emph{Three honest opens survive the deliverable.}

\begin{enumerate}
\item[(F-G2-1)] Higher-genus G2 at $g \ge 3$. The universal KZB
   extends to $\Mbar_{2,n}$ (Enriquez 2013) and to $\Mbar_{g,n}$
   for $g \ge 3$ only via the Beilinson-Drinfeld Green's function
   $G_{C_g}(z,w)$ with open regularity at nodal boundary strata (opus 04
   Cycle 4; CLAUDE.md F5 frontier). Path: Jimbo-Miwa-Stokes
   classification of the irregular strata + Hain 2020 mapping-class-group
   action.

\item[(F-G2-2)] $F_5$ canonical lift at $g = 1$ for class
   $\mathsf{L}, \mathsf{C}, \mathsf{M}$. The obstruction
   $\mathrm{obs}^{(1)}_{\mathrm{double}}$ (Theorem~\ref{thm:f5-obstruction-explicit})
   is non-zero for $r_{\max} \ge 3$, and no canonical elliptic
   trivialisation is known. Path: Ribet 1976 + Mazur-Wiles 1984 for the
   $p$-adic trivialisation at Eisenstein congruences; this
   trivialises the cocycle mod $p$ but not canonically over $\ZZ$.
   Conjecturally the full lift exists after tensoring with the
   Ramanujan discriminant $\Delta(\tau)$ (weight 12), introducing a
   $\wp(\tau)/\Delta(\tau)$ correction that absorbs the elliptic 2-cocycle.

\item[(F-GRT-Eis)] Explicit computation of the Eisenstein phantom gauge
   at primes beyond $691, 3617$. The Kummer-Bernoulli list is
   $\{691, 3617, 43867, 283\cdot617, 131\cdot593, 657931, \ldots\}$
   (next irregular primes at $(B_{18}, B_{20}, B_{22}, B_{24}, B_{26}, B_{28}$);
   OEIS A000928); the shadow tower through $r \le 11$ sees only the
   first two. Higher-$r$ shadow entries would probe the remaining
   primes.
\end{enumerate}

\subsection*{Hidden structure discovered during the attack}

\emph{HS-1: the shadow-tower classes and the level-$n$ Arnold-KZ
matching.}
The shadow-tower index $r_{\max}(A)$ controls the level $n$ at which
G2 first sees new coherence (Theorem~\ref{thm:level-by-level-g2}
structural reading): $r_{\max} = 2$ only sees $n = 2$; $r_{\max} = 3$
sees up to CYBE at $n = 3$; $r_{\max} = 4$ sees the pentagon+hexagon at
$n = 4$; $r_{\max} = \infty$ extends past the Kohno Lie algebra
three-form coherence envelope and requires iterated-integral
coefficients. This is a precise link between the shadow-tower
classification and the per-level Arnold-KZ content of G2.

\emph{HS-2: class $\mathsf{G}$ is the unique $F_5$-globalisable family
past genus zero.}
Theorem~\ref{thm:f5-obstruction-explicit} identifies class $\mathsf{G}$
(Heisenberg) as the only standard-landscape class for which $F_5$
extends canonically to $g = 1$. This is new: the master synthesis
labelled $F_5$ globally OPEN past genus zero without family refinement.

\emph{HS-3: the refined face count of two orbits, and the motivic
bridge.}
Cycle 2 Theorem~\ref{thm:refined-face-count} refines ``seven'' to
``two $\GRTfin$-orbits + one motivic joining element.'' The motivic
joining element is the Brown associator at binary degree; this
identifies F1 (bar hub) and F4 (KZ flat sections) as the \emph{two
primary faces}, and the remaining five as $\GRTfin$-images of these
two. The structural count is: \emph{two} at the level of motivic
primitives, \emph{seven} at the level of algebraically distinct
presentations, \emph{nine} at the level of the full Vol II motivic
+ operadic extension.

\emph{HS-4: the Kummer-Bernoulli primes as Eisenstein phantom gauges.}
Cycle 5 Theorem~\ref{thm:torsor-rigour-eis} identifies the primes
$691, 3617, \ldots$ as the loci of non-faithful action on
$\mathrm{Face}(A)$, inscribing a new bridge between the Vol I
shadow-tower arithmetic and the GRT motivic torsor. The bridge is
\emph{explicit} via the Tamarkin formality-gauge realisation:
$\GRT(\QQ)$ acts on $\mathrm{Face}(A)$ via its action on the
$E_2$-operad deformation moduli, and the $p$-adic specialisation
inherits the Eisenstein-cocycle kernel from Mazur-Wiles.

\emph{HS-5: the critical-level dichotomy in F3.}
Theorem~\ref{thm:yangian-km-nc}(3) states that at critical level
$k = -\hv$, F3 Yangian rational degenerates to the
Feigin-Frenkel-Reshetikhin Bethe subalgebra on
$\mathfrak{z}(\widehat{\fg}) \simeq \mathcal{O}(\mathrm{Op}_\fg(D^\times))$.
The dichotomy is a new entry in the seven-face landscape: at
non-critical level, F3 and F7 are $\GRTfin$-equivalent; at critical
level they are \emph{different mathematical objects}, the former
commutative (Bethe), the latter flat-connection (trace-form without
level prefix). This corresponds, in the Beilinson-Drinfeld programme,
to the passage from the chiral Yangian at non-critical level to the
Langlands dual oper algebra at critical level.

\subsection*{Verification against \texttt{landscape\_census.tex}}

Per-family $r(z)$ check (level prefix mandatory):

\begin{itemize}
\item \textbf{Class $\mathsf{G}$ (Heisenberg, $\cH_k$):}
   $r^{\mathrm{Heis}}(z) = k/z$. Match: line 655 of
   \texttt{landscape\_census.tex}. F1 = F7 (bar residue and Wick
   contraction are literally the same for abelian Heisenberg); F3
   Yangian is trivial (Heisenberg has no simple roots); F5 globalises
   at $g = 1$ (Theorem~\ref{thm:f5-obstruction-explicit}).

\item \textbf{Class $\mathsf{L}$ (affine Kac-Moody,
   $\widehat{\fg}_k$):}
   $r^{\mathrm{KM}}(z) = k\,\Omtr/z$ (trace form, BD 2004 convention)
   $= \Omega/((k+\hv)z)$ (Killing form, KZ convention). Match:
   line 685 (class-L header) and line 689 of
   \texttt{landscape\_census.tex}. F3 Yangian at
   non-critical level: $R(z) = 1 + \hbar\,\OmKill/z + O(\hbar^2)$
   with $\hbar = 1/(k+\hv)$. Critical-level dichotomy per HS-5.

\item \textbf{Class $\mathsf{C}$ ($\beta\gamma$, $bc$ ghosts):}
   $r^{\beta\gamma}(z) = 0$ (regular; symplectic normalisation absorbs
   level). Match: line 670--672 of \texttt{landscape\_census.tex}. F3
   Yangian trivial; F5 globalises at $g \le 1$ only for diagonal
   $\beta\gamma$ (symplectic form constant); obstructed for
   $\beta\gamma(\lambda)$ with $\lambda \ne 1/2$.

\item \textbf{Class $\mathsf{M}$ (Virasoro, $\mathrm{Vir}_c$):}
   $r^{\mathrm{Vir}}(z) = (c/2)/z^3 + 2T/z$. Match: lines 699--701 of
   \texttt{landscape\_census.tex}. Odd-pole parity (all odd poles)
   from bosonic-parity constraint (Remark
   \ref{rem:bosonic-parity-census} of landscape\_census). F3 quantum
   group: the Virasoro deformation at generic central charge has no
   rational Yangian form; the appropriate quantisation is the
   \emph{q-deformed Virasoro} of Shiraishi-Kubo-Awata-Odake 1996
   (q-Virasoro), which sits in a different column of the associator
   torsor. F5 obstructed for generic $c$; globalises at $c = 0$
   (conformal anomaly rigidity) and at $c = -22/5$ (non-unitary
   minimal Yang-Lee).

\item \textbf{Class $\mathsf{W}$ ($\mathcal{W}_N$, $\mathcal{W}_3$
   $WW$ channel):}
   $r^{\mathcal{W}_3, WW}(z) = (c/3)/z^5 + 2T/z^3 + \partial T/z^2 + \ldots$
   Match: lines 714--716 of \texttt{landscape\_census.tex}. Even-order
   pole at $z^{-2}$ from $\partial T/(z-w)^3$ term; cf.\ Remark
   \ref{rem:ww-even-poles-census}. F3 Yangian: the quantum W-algebra of
   Frenkel-Reshetikhin 1998 is the rational form, with $R(z)$-matrix
   living in $U_q(\widehat{\mathcal{W}}_N)$; critical-level dichotomy
   per HS-5 via the Feigin-Frenkel duality
   $\mathcal{W}_k(\fg) \simeq \mathcal{W}_{k^\vee}(\fg^\vee)$
   relating non-critical levels in dual pairs.

\item \textbf{Yangian $Y(\fg)$:}
   Row 562 of \texttt{landscape\_census.tex} records
   $Y(\fg)$ as class $\mathsf{M}$ with $r_{\max} = \infty$ and
   ``$R$-matrix $= \Theta_\cA|_{E_1}$.'' The F3 presentation
   matches the row directly; the F1-F3 identification is via the
   quasi-classical limit $\hbar = 1/(k+\hv) \to 0$.
\end{itemize}

All per-family entries verified against
\texttt{chapters/examples/landscape\_census.tex}, level prefixes
present and canonical.

\subsection*{Epilogue}

The two threads of the master synthesis (G2 Arnold-KZ identity; seven
faces of $r(z)$ as a $\GRT(\QQ)$-torsor) meet in a single universal
element $r(z) = \Res^{\coll}_{0,2}(\Theta_A)$, whose bar-differential
incarnation $\dbar = \KZ^{*}(\nabla_{\Arnold})$ is strict at chain
level (Theorem~\ref{thm:g2-healed-chain-local}), globally extended
past genus $2$ via the universal KZB of Felder-Wieczerkowski + CEE
(Theorem~\ref{thm:g2-healed-infty-global}), and whose seven-face
presentation is a two-orbit $\GRTfin$-gauge equivalence at chain level
with a $\Gmot_{\MT(\ZZ)}/\mathrm{Eis}_p$ phantom gauge at irregular
primes (Theorem~\ref{thm:seven-faces-two-scope},
Theorem~\ref{thm:torsor-rigour-eis}).

Three frontiers (F-G2-1, F-G2-2, F-GRT-Eis) remain honest opens. Five
hidden structures (HS-1 through HS-5) emerge from the attack-heal
cycle, each of independent interest: the shadow-tower / Arnold-KZ
per-level bridge, the class-$\mathsf{G}$ $F_5$-globalisation dichotomy,
the two-orbit motivic refinement, the Kummer-Bernoulli phantom gauge,
and the critical-level Yangian / Bethe dichotomy.

The seven faces remain seven. The count is literary, not
structural, and the structural count is two orbits plus a motivic
bridge. The GRT torsor action is rigorous at every prime, with explicit
phantom gauges at Kummer-Bernoulli primes. The pullback identity is a
tower of theorems, one per level $n$, exhausted by pentagon+hexagon at
$n = 4$. All five attack surfaces survive the heal intact; the healed
statements sharpen, do not retreat from, the master synthesis.

%% ====================================================================
\subsection*{Cycle log (raw)}
%% ====================================================================

\begin{enumerate}
\item[C1] Level-$n$ matching: healed via Theorem~\ref{thm:level-by-level-g2}.
   Three levels ($n=2, 3, 4$) exhaust new content; higher levels
   coherence-forced. Opus 04 Cycle 1 provides the $n=2$ base;
   Cycle 1 here extends through hexagon + pentagon.

\item[C2] Seven-face count: refined to two $\GRTfin$-orbits
   (Theorem~\ref{thm:refined-face-count}). Historical seven preserved
   for reader orientation; Vol II nine-face extension acknowledged;
   structural two-count is the honest answer.

\item[C3] Yangian F3 vs trace-form F7: reconciled via Casimir rescaling
   + Sugawara level shift at non-critical level
   (Theorem~\ref{thm:yangian-km-nc}); critical-level dichotomy
   identified as Feigin-Frenkel-Reshetikhin Bethe subalgebra (new
   hidden structure HS-5).

\item[C4] $F_5$ at $g \ge 1$: explicit cohomological obstruction
   $\mathrm{obs}^{(1)}_{\mathrm{double}} \in H^2(\grt^{\mathrm{ell}}, \mathfrak{sp} \otimes \mathfrak{sp})$
   (Theorem~\ref{thm:f5-obstruction-explicit}). Class $\mathsf{G}$ is
   the unique globalisable family past genus zero (hidden structure
   HS-2). Higher genus honest open (F-G2-2).

\item[C5] GRT torsor rigour + Kummer-Bernoulli primes: torsor rigorous
   at every scope; exotic behaviour at $p = 691, 3617$ is explicit
   Eisenstein phantom gauge (Theorem~\ref{thm:torsor-rigour-eis}). Vol
   I shadow tower through $r = 11$ witnesses the first two irregular
   primes; higher-$r$ shadow entries probe the remaining irregular
   list (F-GRT-Eis).
\end{enumerate}

End of deliverable.
